\chapter*{Glosario}
\addcontentsline{toc}{chapter}{Glosario}
\markboth{Glosario}{Glosario}

Los términos siguientes se utilizan con este significado en la memoria y en los anexos.

\begin{longtable}{p{0.24\textwidth}p{0.65\textwidth}}
\toprule
\textbf{Término} & \textbf{Significado en el proyecto}\\
\midrule
\endfirsthead
\toprule
\textbf{Término} & \textbf{Significado en el proyecto}\\
\midrule
\endhead
\bottomrule
\endlastfoot
Bundle & Colección persistida de aplicaciones, oficial o propiedad de una cuenta; se distingue de la selección temporal de una descarga.\\
Candidato & Enlace o alternativa de instalación detectada para una aplicación antes de su validación final.\\
Enlace mágico & Enlace de acceso enviado por correo, de un solo uso y vigencia limitada; su consumo crea o autentica una cuenta de usuario.\\
Fallback literal & Repetición completa de una búsqueda en modo literal cuando la ruta semántica no puede ofrecer un resultado coherente.\\
Fuente & Página o servicio del editor y configuración asociada desde los que se intenta obtener un instalador.\\
Índice semántico & Proyección reconstruible de documentos y vectores que permite proponer candidatos a una búsqueda; MySQL conserva los filtros y estados públicos.\\
Riesgo residual & Riesgo estimado que permanece después de aplicar una medida; no equivale por sí solo a una prueba superada.\\
Trabajo & Solicitud asíncrona de descarga con estados, elementos y un resultado temporal propio.\\
Trazabilidad & Relación verificable entre requisitos, casos de uso, riesgos, pruebas, resultados y evidencias.\\
\end{longtable}
